\documentclass[11pt]{article}
\usepackage[utf8]{inputenc}
\usepackage[T1]{fontenc}
\usepackage{lmodern}
\usepackage{textcomp}
\usepackage{amsmath, amssymb, amsthm}
\usepackage{graphicx}
\usepackage{listings}
\usepackage[hidelinks]{hyperref}

\theoremstyle{plain}
\newtheorem{theorem}{Teorema}[section]
\newtheorem{lemma}[theorem]{Lema}
\newtheorem{proposition}[theorem]{Proposición}
\newtheorem{corollary}[theorem]{Corolario}

\theoremstyle{definition}
\newtheorem{definition}[theorem]{Definición}


\title{c9}
\date{}
\begin{document}
\maketitle

\section{Chapter 9}

\section{Existence and Definitions}

In this chapter we take a closer look at existence statements. These statements assert that there exists a quantity x that has a certain property P. In symbols:

$\exists$x $\in$X, P(x)

where X is a set and P is a predicate over X. For example, the statement

The integer 109 is the sum of two primes

asserts the existence of a pair of prime numbers with a given sum (what are X and P in this case?).

We shall discuss proofs of existence and the connection between existence and definitions.

\section{9.1 Proofs of Existence}

To prove that something exists it is not necessary that the object in question be constructed explicitly. For example, Euclid's proof of the infinitude of primes (Sect. 7.5) establishes the existence of infinitely many things without producing any of them. By contrast, we proved that Euler's polynomial n2 + n + 41 is composite for infinitely many values of n by exhibiting such values explicitly (Eq. 7.8).

Accordingly, existence proofs are said to be constructive if an explicit construction is given, and non-constructive if no explicit construction is given. The constructive method may be described as a two-stage process:

(WHAT?) Identify an element x of X. (WHY?) Show that P(x) is True.

The format of the proof should make clear which part is the WHAT and which part is the WHY.

© Springer-Verlag London 2014 F. Vivaldi, Mathematical Writing, Springer Undergraduate Mathematics Series, DOI 10.1007/978-1-4471-6527-9\_9

151

152 9 Existence and Definitions

We begin with a constructive existence proof [31, p. 13]:

Theorem. There is a real number $\alpha$ such that $\alpha$2 = 2.

\begin{proof}
Proof. Draw a square of side 1.
\end{proof}

(WHAT?) Let $\alpha$ be the length of a diagonal of the square. (WHY?) Then by Pythagoras, $\alpha$2 = 2. 

This minimalist proof takes for granted that the length of a diagonal of the unit square is a real number. A justification of this step requires the tools of analysis.

Let us re-visit the counterexample (7.8) concerning Euler's polynomial.

Theorem. There are infinitely many integers n for which p(n) = n2 + n + 41 is composite.

\begin{proof}
Proof.
\end{proof}

(WHAT?) Let n = 41k, for k $\in$N. (WHY?) Then p(41k) = 41(41k2 + k + 1), not a prime. 

In this proof the WHAT part identifies an infinite sequence of integers.

Theorem. If two integers are the sum of two squares, then so is their product.

This statement is an implication; the existence statement in the deduction is conditional to another existence statement in the hypothesis. We give a concise constructive proof, where an eloquent formula supplies at once the WHATs and the WHYs.

\begin{proof}
Proof. This statement follows from the algebraic identity
\end{proof}

( j2 + k2)(m2 + n2) = ( jm + kn)2 + (km $-$jn)2. 

Many non-constructive existence proofs employ contradiction. One assumes that the object being defined does not exist, and derives a false statement. It is not surprising that an argument of this kind could fail to provide information about the object in question.

However, not all non-constructive existence proofs use contradiction.

$\surd$

2 is rational.

Theorem. There is an irrational number a such that a

$\surd$

$\surd$

2. We have two cases:

\begin{proof}
Proof. Consider
\end{proof}

2

$\surd$

$\surd$

2 is rational. Then a = $\surd$

CASE I:

2

2 is as required.

$\surd$

$\surd$

2 is irrational. Then

CASE II:

2

2$\surd$ 2

$\surd$

$\surd$

$\surd$

2 = ( $\surd$

$\surd$

$\surd$ 2

2)2 = 2,

= (

2)

2

$\surd$

$\surd$

2 is as required. 

so a =

2

9.1 Proofs of Existence 153

This is a proof by cases---see Sect. 7.2. We note that we aren't given a WHAT. There are two different WHATs, and we have no idea which of them works for the WHY.

Our next example is a well-known existence statement, Dirichlet's box principle (or the pigeon-hole principle).

Theorem. Given n boxes and m objects in them, if m > n, then at least one box contains more than one object.

The proof of this rather self-evident implication is necessarily non-constructive. We use contradiction, namely the both ends method described in Sect. 7.5. It's easy to predict where contradiction will lead us.

\begin{proof}
Proof. Let m and n be integers, and let m j be the number of objects in the jth box. Assume that m > n but m j $\leq$1, for j = 1, . . . , n. Then
\end{proof}

n

n

m =

m j $\leq$

1 = n

j=1

j=1

giving m $\leq$n, contrary to the assumption. 

In this proof we are not given the WHAT, and hence there cannot be a WHY step either. The proof begins by introducing the symbol m j; this simple but important step makes the rest immediate.

The non-constructive proof of the following statement employs Dirichlet's box principle.

Theorem. In any set of integers with more than n elements there must be two integers whose difference is divisible by n.

\begin{proof}
Proof. Let n be given, and let S be a set of m integers, with n < m. We divide each element of S by n, obtaining m integer remainders. These remainders can assume at most n distinct values, and therefore two elements of S, say xi and x j, must have the same remainder, by Dirichlet's box principle. But then xi $-$x j gives reminder zero when divided by n, as desired. 
\end{proof}

Solving equations is a quintessential mathematical task. In some cases we may be looking for a particular number, which means that we are after a constructive existence proof. In other cases we may be satisfied by a proof that a solution exists at all, or that it exists in a specified range of values of the argument. Let us examine a constructive and a non-constructive existence proof of the same statement.

Proposition. The equation x4 $-$10x2 + 1 = 0 has four distinct real solutions.

Constructive Proof. Let f (x) = x4 $-$10x2 + 1. (WHAT?) Let x$\pm$ =

$\surd$ 3 $\pm$

$\surd$

2; then x$\pm$ are real numbers. (WHY?) We compute, using the binomial theorem

154 9 Existence and Definitions

$\surd$

$\surd$ 3 $\pm$

2)4 $-$10( $\surd$

$\surd$ 3 $\pm$

2)2 + 1

f (x$\pm$) = (

$\surd$

$\surd$ 6 + 36 $\pm$ 8

$\surd$ 6 + 4 $-$30 $\mp$20

= 9 $\pm$ 12

6 $-$20 + 1

= 0

where the top and bottom signs in $\pm$ and $\mp$match. The above calculation shows that x$\pm$ are roots of f . Since the real numbers x$\pm$ are positive and distinct, and f is an even function, we conclude that $-$x$\pm$ are also roots of f . 

This constructive existence proof involved `guessing' the solutions, and then verifying that the guess was correct.1

Non-Constructive Proof. Let f (x) = x4 $-$10x2 + 1. The computation

f (0) = 1 f (1) = $-$8 f (4) = 97,

shows that f changes sign twice in the interval (0, 4). Because f is continuous, there exist two distinct real numbers x$-$, x+ in the open intervals (0, 1) and (1, 4), respectively, at which the function f vanishes. The real numbers $-$x$\pm$ are also roots of f, because f is even. The four numbers $\pm$x$\pm$ are clearly distinct. 

The key argument in the proof was inferring the existence of a root from a change of sign of a continuous function. This is the intermediate value theorem, a non-constructive existence theorem of real analysis [26, Theorem 4.4]:

Theorem. Anyreal-valuedcontinuousfunction f onaclosedinterval[a, b]assumes every value between f (a) and f (b).

Our non-constructive proof provides some information about the solutions, in the form of bounds. These bounds can be sharpened by doing some extra work; e.g., with f as above we have

7

6

= $-$503

1657 =

f

f 234256

22

19

130321

and hence

19 < x$-$< 7 6 22.

Considering that 7/22 $-$6/19 < 3 $\times$ 10$-$3, our knowledge of the solution x$-$has increased markedly. We see that the distinction between a constructive and a non-constructive proof is blurred.

A non-constructive existence statement that provides some information about an object---typically a number---in the form of bounds, is said to be effective.2 The

1 There was no guess here: first I chose x+, and then I derived the equation of which x+ is a solution. 2 Some authors use `effective' to mean `constructive'.

9.1 Proofs of Existence 155

above non-constructive proof is effective. Likewise, the estimate of the size of the nth prime number given by Theorem 8.1 (Sect. 8.4) may be characterised as an effective version of Euclid's theorem.

In the mathematics literature the meaning of `constructive proof' is sometimes extended to include methods that, while not providing an actual object, provide an algorithm for constructing the object.

\section{9.2 Unique Existence}

Unique existence is a stronger version of existence:

There is exactly one element of X having property P.

where, as usual, X is a set and P is a predicate over X. For example:

The equation f (x) = 0 has exactly one solution. The sets A and B have only one point in common. The function g has precisely one stationary point. The matrix M has a unique real eigenvalue greater than 1.

The expressions `exactly', `only', `precisely' differentiate unique existence from existence.

Unique existence is a hidden conjunction; the conjuncts are

Existence: There is one element of X with property P. Uniqueness: If x, y $\in$X have property P, then x = y.

The corresponding symbolic expression is overloaded:

($\exists$x $\in$X, P(x)) $\land$($\forall$x, y $\in$X, (P(x) $\land$P(y)) $\Rightarrow$(x = y)) (9.1)

and to simplify it we introduce a new quantifier $\exists$! denoting unique existence:

$\exists$! x $\in$X, P(x).

For example, given a function f : X $\to$Y, the concise expression

$\forall$y $\in$Y, $\exists$! x $\in$X, f (x) = y

states that f is bijective.

In Chap. 8 we proved (twice) that any integer n greater than 1 is a product of primes. For each n, this statement asserts the existence of an unspecified number of primes. The fundamental theorem of arithmetic upgrades existence to unique existence.

156 9 Existence and Definitions

Theorem. Every natural number greater than 1 may be written as a product of prime numbers. This representation is unique, apart from re-arrangements of the factors.

The expression `apart from re-arrangements of the factors' gives the impression of lack of uniqueness. This is not the case: the theorem states that there exists a unique multiset of prime numbers (see Sect. 2.1) with the property that the product of its elements is the given natural number. Alternatively, there is a unique set of pairwise co-prime prime-powers with the same property.

In a proof of unique existence the conjuncts are normally proved separately, in either order.

Theorem. The identity element of a group is unique.

The group axioms were given in Sect. 7.1. The existence of an identity element does not require a proof, being an axiom (axiom G3). We only have to prove uniqueness.

\begin{proof}
Proof. Let us assume that a group has two identity elements, 1 and 2. Applying the axiom G3 to each identity element, we obtain
\end{proof}

1 2 = 1 1 2 = 2

and hence 1 = 2. 

Theorem. For every set X and subset Y, there is a unique set Z such that Y $\cup$Z = X and Y $\cap$Z = $\emptyset$.

\begin{proof}
Proof. Let Z = X$\setminus$Y. By construction, Y $\cap$Z = $\emptyset$. Since Y $\cup$Z is a subset of X, and every element of X belongs to either Y or Z, we have Y $\cup$Z = X. So a set Z with the required properties exists.
\end{proof}

To prove uniqueness, assume that Z is a set with the stated properties. Then, since Y $\cup$Z = X, Z must be a subset of X, and it must contain all elements of X that aren't in Y. So Z must contain X$\setminus$Y.

Since Y $\cap$Z = $\emptyset$, the set Z can't contain any elements of X that are in Y. So Z must be X$\setminus$Y. 

The existence part of this proof is constructive. For this reason, the proof of uniqueness does not establish the second conjunct in (9.1) explicitly. Rather than considering two objects with the stated properties and showing that they are the same, we consider a single object and show that it is the same as the one defined in the existence part.

Theorem (Fermat). The only natural numbers x for which 1 + x + x2 + x3 is a square are x = 1 and x = 7.

Proving existence is immediate: 1 + 1 + 12 + 13 = 22, 1 + 7 + 72 + 73 = 202. Uniqueness is much harder, see [12, pp. 38, 382].

9.3 Definitions 157

\section{9.3 Definitions}

Definitions are closely related to unique existence. When we define a symbol or a name, we must ensure that this quantity actually exists, and that it has the stated properties. A definition may identify a unique quantity (Let a be the length of a diagonal of a regular pentagon with unit area), or specify membership to a unique non-empty set (Let p be a prime such that p + 2 is also prime). If these conditions are met, then our object is well-defined; otherwise it is ill-defined.

Let us begin with definitions of sets. Here there is a distinctive safety net: if the conditions imposed on a set are too restrictive, then this set will be empty rather than ill-defined. Nonetheless, we should be alert to this possibility, because the consequences of an empty definition could be more serious than formally correct nonsense. For example, consider the following definition:

Let M be the set of 2 by 2 integral matrices with odd entries and unit determinant.

The set M is empty because the determinant of a matrix with odd entries is an even integer, and so it cannot be 1. It's easy to run into trouble now: Let M $\in$M .

Sets defined by several conditions are commonplace; for instance, the set of solutions of a system of simultaneous equations is the intersection of the solution sets of the individual equations. A considerate definition should flag the possibility of an empty intersection:

Let A1, . . . , An be sets, and let A denote their (possibly empty) common intersection

A = A1 $\cap$A2 $\cap$\textperiodcentered{} \textperiodcentered{} \textperiodcentered{} $\cap$An. (9.2)

For n > 2, the set A is well-defined, because the intersection operator is associative. For n = 1 the Formula (9.2) is, strictly speaking, unspecified, but the interpretation

A = A1 is quite natural, so a separate treatment is unnecessary. One could even stretch this definition to mean A = $\emptyset$if n = 0, although a clarifying remark would be needed in this case.

Next we consider an innocent-looking integer sequence whose existence is more delicate than expected.

Let qn be the smallest prime number with n decimal digits.

Does such a prime exist for all n? Let's look at the first few terms of the sequence:

(qn)$\infty$

1 = (2, 11, 101, 1009, 10007, . . .). (9.3)

In all cases, the prime qn lies just above 10n$-$1, and it seems unavoidable that there is at least one prime between 10n$-$1 and 10n.

158 9 Existence and Definitions

However, given that there are arbitrarily large gaps between consecutive primes,3

an argument is needed to rule out the possibility that a large gap could include all integers with n digits, for some n. Such an argument is not elementary, and we shouldn't keep the reader pondering on this. So the definition of the sequence (qn) should incorporate a remark or a footnote of the type:

This sequence is well-defined due to Bertrand's postulate4: for all n $\geq$1 there is at least one prime p such that n < p $\leq$2n [15, p. 343].

A function definition f : A $\to$B requires specifying two sets, the domain A and the co-domain B, as well as a rule that associates to every point x $\in$A a unique point of f (x) $\in$B. For the function f to be well-defined, we must ensure that the specification of A and of the rule x $\to$f (x) do not contain any ambiguity.

By contrast, there is flexibility in the specification of the co-domain B, in the sense that any set containing f (A) may serve as a co-domain. Formally, different choices of B correspond to different functions, although such distinctions are often unimportant. But then why don't we always choose f (A) as co-domain? This would have the advantage of making every function surjective. The problem is that we may not know what f (A) is, or the description of f (A) may be exceedingly complicated.

For instance, let us return to the sequence (qn) given in (9.3). Writing

qn = 10n$-$1 + an (9.4)

we find

(an)$\infty$

1 = (1, 1, 1, 9, 7, . . .).

Let us now define the function

f : N $\to$N f (n) = an.

This function is well-defined but we have limited knowledge of the image f (N). It can be shown that f (N) cannot intersect any of the sets 2N, 2 + 3N, and 5N, but we don't know if there are other constraints---see Exercise 10.4.3.

\section{9.4 Recursive Definitions}

A recursive definition is a form of definition connected to the principle of induction. With this method it is possible to define very complex objects---sequences, typically---from minimal ingredients.

3 For n > 1, the integer 1 + n! is followed by n $-$1 composite integers. (Why?) 4 Conjectured by Joseph Bertrand (French: 1822--1900) in 1845. Proved by Pafnuty Chebyshev (Russian: 1821--1894) in 1850.

9.4 Recursive Definitions 159

We have pointed out that a sequence (ak) corresponds to a function over N, whereby ak represents the value of the function at the point k $\in$N. The existence of such a function does not necessarily imply that its values are easily recoverable from the definition of the sequence. For example, we don't know any `useful' way to express the kth prime number as an explicit function of k. In absence of an explicit formula for the elements of a sequence, we seek to represent ak+1 in terms of ak. These are recursive sequences, which are defined by two data: the first term of the sequence, called the initial condition, and the rule which determines a term from the previous one.

For instance, the factorial sequence n! = 1\textperiodcentered{}2\textperiodcentered{}3 \textperiodcentered{} \textperiodcentered{} \textperiodcentered{} n, can be defined recursively as follows:

0! = 1 (k + 1)! = (k + 1) \textperiodcentered{} k! k $\geq$0. (9.5)

This definition, in effect, specifies the order in which the product is computed. Likewise, the nth derivative g(n) of a function g is defined by the recursive rule

g(0)(x) = g(x) g(k+1)(x) = d

dx g(k)(x) k $\geq$0.

More generally, given any set X and any function f : X $\to$X, we define a recursive sequence (ak) of elements of X as follows:

a1 $\in$X given; ak+1 = f (ak), k $\geq$1. (9.6)

The first term of the sequence (the initial condition) is given. Because f is a function, if ak is well-defined for some k $\geq$1, then so is ak+1. The induction principle then ensures that the whole sequence is well-defined. Since distinct initial conditions in

X give rise to distinct sequences, there are as many sequences as there are elements of X: plenty of sequences from just one function!

The term ak+1 of a recursive sequence may depend on d $\geq$1 preceding terms, not just the last one:

ak+1 = f (ak, ak$-$1, . . . , ak$-$d+1). (9.7)

In this case we have f : Xd $\to$X, and d is called the order of the sequence. In a recursive sequence of order d, the first d terms of the sequence must be supplied explicitly because they don't have the required number of preceding terms. These sequences have d initial conditions and they are well-defined by virtue of strong induction.

The Fibonacci sequence5 is a second-order sequence:

a0 = 1; a1 = 1; ak+1 = f (ak, ak$-$1) = ak + ak$-$1, k $\geq$1.

5 Leonardo Pisano, known as Fibonacci (Italian: 1170--ca.1240).

160 9 Existence and Definitions

We find

(1, 1, 2, 3, 5, 8, 13, 21, 34, 55, 89, 144, . . .).

Because the term ak$-$1 = ak+1 $-$ak is a well-defined function of the following two terms, we can extend the Fibonacci sequence backwards to obtain a doubly-infinite sequence:

(. . . , $-$21, 13, $-$8, 5, $-$3, 2, $-$1, 1, 0, 1, 1, 2, 3, 5, 8, 13, 21, . . .).

A recursive sequence that can be extended backwards is said to be invertible. Thus the Fibonacci sequence is invertible.

Recursive sequences may be quite unpredictable, as the following example illustrates. Let

x/2 if x is even 3x + 1 if x is odd.

f : N $\to$N x $\to$

We consider the first-order recursive sequence associated with this function:

x0 = n xt+1 = f (xt), t $\geq$0. (9.8)

The initial condition x0 = 1 leads to a periodic sequence

(1, 4, 2, 1, 4, 2, 1, . . .) (9.9)

consisting of indefinite repetitions of the pattern (1, 4, 2). If instead we start with x0 = 7, we find

(7, 22, 11, 34, 17, 52, 26, 13, 40, 20, 10, 5, 16, 8, 4, 2, 1, 4, 2, 1, . . .)

so, after an irregular initial excursion, the sequence settles down to the same periodic pattern as the previous sequence. The analysis of the sequences (9.8) for general initial condition n presents formidable difficulties, which are distilled into a famous unsolved problem, the so-called `3x+1' conjecture [43]:

Conjecture. For any n, the sequence (9.8) contains 1.

This conjecture states that all these sequences are eventually periodic and that they reach the same periodic pattern.

We note that any open conjecture leads to objects whose existence is at present undecidable, see Exercises 9.4 and 10.4.

9.5 Wrong Definitions 161

\section{9.5 Wrong Definitions}

AswedidinSect.7.7forlogicalarguments,weconsiderheresomecommonmistakes made in definitions. A faulty definition may imply that the object being defined does not exist at all, or that it is not the one we had in mind, or that there is more than one object that fits the description.

We begin with an incorrect function definition which can be put right in several ways.

Wrong Definition. Let f be given by:

f : Q $\to$Q f (x) = x2y + 1.

The mistake is plain but fatal: the expression f (x) involves the unspecified quantity y. As defined, f is not a function. There are many legitimate interpretations of what this formula could mean, each resulting in a different function.

1. We re-define the domain, supplying the missing information via a second argu-

ment.

f : Q2 $\to$Q f (x, y) = x2y + 1.

In this definition y is also rational; clearly there are other possibilities.

2. The missing argument is regarded as a parameter (see Sect. 6.1).

f$\lambda$: Q $\to$Q f$\lambda$(x) = $\lambda$x2 + 1 $\lambda$ $\in$Q.

We have highlighted the change in status of the variable y by turning it into a subscript and switching to the Greek alphabet. For every value of $\lambda$, we have a different function.

3. The missing argument is regarded as an indeterminate. In this case f (x) is a polynomial in y, and we must re-define the co-domain of f accordingly.

f : Q $\to$Q[y] f (r) = r2y + 1.

The symbol Q[y] denotes the set of all polynomials in the indeterminate y with rational coefficients (see Sect. 3.1). To avoid confusion, we have replaced x with a symbol that normally is not used for an indeterminate, and which reminds us that this variable is rational.

In the second example, the existence of a function requires more conditions than those given. The definition relies on a hidden assumption, so that using it adds extra information (we say that the definition is creative).

Wrong Definition. Let A and B be sets, and let f, g : A $\to$B be functions. We define the function h = f + g as follows:

162 9 Existence and Definitions

h: A $\to$B x $\to$f (x) + g(x).

The hidden assumption is that the elements of the co-domain can be added together. However, B may be a set where addition is not defined (e.g., f and g are predicates) or which is not closed under addition (e.g., B is an interval). We give a more restrictive definition without hidden assumptions.

Definition. Let A be a set, let B be a set closed under addition, and let f, g:

A $\to$B be functions. \dots{}

In the next example we attempt to extend to modular arithmetic the concept of the reciprocal of an element, but our definition has a hidden assumption.

Wrong Definition. Let m be a natural number, and let a be an integer. If [a]m = [0]m, then we define the multiplicative inverse [a]$-$1

m as follows:

def = [b]m where [a]m[b]m = [1]m.

[a]$-$1

m

Take m = 6 and a = 3. We have [3]6 = [0]6, but the equation [3]6[x]6 = [1]6 has no solution. What went wrong? We have assumed that the implication

if a = 0 and ab = ac, then b = c,

which is valid for real or complex numbers, is also valid for congruence classes. Requiring that [a]m = [0]m, namely that a is not divisible by the modulus is not enough; for this implication to hold, we must assume that a is co-prime to the modulus.

Definition. Let m be a natural number and let a be an integer co-prime to m. We define the multiplicative inverse [a]$-$1

m of [a]m as follows\dots{}

In the following definition both existence and uniqueness are problematic.

Wrong Definition. Let p $\in$Z[x], and let z p be the root of p having smallest modulus.

If p has degree zero, then z p does not exist, so this definition relies on the hidden assumption that p has positive degree. If z p exists, then it may still not be unique (e.g., p = x2 + 1), but the use of the definite article (`the root') implies that it is. Requiring that p be non-constant solves the existence problem. The way we address uniqueness depends on the context. If we merely require that no root of p has smaller modulus than z p, then we don't need uniqueness.

Definition. Let p $\in$Z[x] with deg(p) > 0, and let z p be a root of p with smallest modulus.

This definition does not identify z p; it tells us that z p is a member of a well-defined non-empty set, but this set may have more than one element. If, on the other hand, we require a specific complex number, then we must supply additional information.

9.5 Wrong Definitions 163

In the following definition we impose a condition on the argument of z p, which removes any ambiguity.

Definition. Let p be a non-constant polynomial with integer coefficients, and let z p be the root of p with smallest modulus. If there is more than one root with this property, we let z p be the root with smallest non-negative argument.

In our final example, the definition hides a subtle lack of uniqueness.

Wrong Definition. Let x $\in$[0, 1], and let f be the function that associates to x its binary digits sequence

f : [0, 1] $\to$\{0, 1\}N x $\to$(c1, c2, . . .),

where

$\infty$

ck

x =

ck $\in$\{0, 1\}. (9.10) 2k

k=1

(The co-domain of f is the set of all infinite binary sequences, see Sect. 3.5.1.) What's wrong with this definition? Consider the rational numbers

$\infty$

1

xn = 1

2n =

n $\geq$1. 2k+n

k=1

If we write xn in the form (9.10), then the above identity shows that there are two binary representations, namely

, 1, 0, 0, 0, . . .) and (0, . . . , 0

, 0, 1, 1, 1, . . .). (9.11)

(0, . . . , 0

n$-$1

n$-$1

So the function f is not uniquely defined at these rationals and, more generally, at all rationals of the form

n$-$1

ck 2k + xn

k=1

where c1, . . . , cn$-$1 are arbitrary binary digits.

We resolve this ambiguity by choosing consistently the first of the two digit sequences in (9.11).

Definition. Let x $\in$[0, 1) be given by

$\infty$

ck

x =

ck $\in$\{0, 1\}. 2k

k=1

Without loss of generality, we assume that all sequences (ck) contain infinitely many zeros. Let now f be the function that associates to x its binary digit sequence

164 9 Existence and Definitions

f : [0, 1) $\to$\{0, 1\}N x $\to$(c1, c2, . . .).

Exercise 9.1 Use Dirichlet's box principle to prove the following existence statements.

1. Given any five points in the unit square, there are two points whose distance apart

$\surd$

is at most 1/

2.

2. In any finite group of people, there are two people with the same number of

friends.

3. Every recursive sequence of elements of a finite set is eventually periodic.

Exercise 9.2 Each function definition contains an error. Explain what it is and how it should be corrected.

1

1. f : R $\to$R x $\to$ x2 + x $-$1

x2 $-$1

2. f : R $\to$R x $\to$

a

3. f : Q $\to$Z b $\to$|a $-$b|

n

$\lfloor$kx$\rfloor$

4. f : Z $\times$ R $\to$Z (x, n) $\to$

k=1

5. f : Z[x] $\times$ N $\to$N (p, n) $\to$gcd(p, xn)

6. f : Z2 $\to$Z (m, n) $\to$nZ $\cup$mZ.

Exercise 9.3 The following definitions have several flaws. Explain what they are, hence write a correct, clearer definition.

1. Let a = (b1, b2, . . .) be a given sequence, with bk $\in$B. We define the function

m

b$-$1

f : Z $\to$B f (m) =

k .

k=1

2. Let X be a subset of Q, and let f (X) be the number of integers in X, namely

f : Q $\to$N f (X) = \#(x $\in$X $\cap$x $\in$Z).

3. Let l be a line in the cartesian plane, let P, Q $\in$R2 be the points of intersection of l with the coordinate axes, and let F(l): R2 $\to$Rbe the function that gives the length of the segment joining P and Q.

Exercise 9.4 Exploit each conjecture to define a function whose existence cannot be decided at present.

1. The twin-primes conjecture (Sect. 7.6).

2. The Goldbach conjecture (Sect. 7.6).

3. The 3x + 1 conjecture (Sect. 9.4).

\end{document}